\section{附录}

\subsection{伪标签 CSV 的数据契约}

\texttt{process\_esc\_csr.py} 中的 \texttt{get\_sent\_ground} 对格式有硬性要求，
任一字段不符都会触发静默回退。列定义见表~\ref{tab:csv}。

\begin{table}[H]
\centering
\caption{伪标签 CSV 的列定义}
\label{tab:csv}
\small
\begin{tabular}{cll}
\toprule
列号 & 列名 & 要求\\
\midrule
0 & \texttt{id} & 行号，无实义\\
1 & \texttt{doc\_id} & \textbf{被读取}，须与 pkl 中的 \texttt{doc\_id} 完全一致\\
2 & \texttt{text} & 原文，不参与计算\\
3 & \texttt{sentences} & 句子列表\\
4 & \texttt{event\_pair} & 事件对，不参与计算\\
5 & \texttt{node\_event\_id} & 事件编号，不参与计算\\
6 & \texttt{answer} & \textbf{被读取}，须为 \texttt{str(list[str])} 形式\\
\bottomrule
\end{tabular}
\end{table}

第 6 列的格式最容易出错：它必须是一整篇文档所有事件对答案的列表的字符串表示，
例如 \texttt{"['1,2', 'None', '3']"}，其中每个元素是一对事件的答案——
正例给句子编号（逗号分隔），负例给字符串 \texttt{'None'}。
一条 \texttt{doc\_id} 对应一行，\texttt{answer} 列表的长度必须等于该文档的事件对数，
否则标签矩阵与事件对数不匹配，训练会直接报错或（更糟）静默错位。

\subsection{LLM 调用的提示词与响应样例}

以下为 \texttt{raw\_dev\_sen.jsonl} 中的一条真实记录（文档 \texttt{37\_1}，
事件对 \texttt{[3,1]}，模型 qwen3.8-27b，共 27 句）。

\paragraph{系统提示}
\begin{lstlisting}
You are an expert with few words and will only answer questions accurately
\end{lstlisting}

\paragraph{用户提示（节选）}
\begin{lstlisting}
The input document is labeled with an ordinal number before each sentence,
e.g. (0,1...) , knowing that event: earthquake and event: injured in the 2
and 0 sentences are causally related, earthquake means [A powerful earthquake
in Indonesia's Aceh province flattened buildings and sparked landslides
Tuesday, killing at least five people and injuring dozens in a region
devastated by the quake-triggered tsunami of 2004.]. injured means[Five dead,
dozens injured ...]. identify which sentences from the document articulate
the causal relationship between events earthquake and injured. ...
\end{lstlisting}

\paragraph{原始响应}
\begin{lstlisting}
[0, 2]
\end{lstlisting}

两个事件分别位于第 2 句与第 0 句，LLM 返回的集合恰好是 $\{0,2\}$——
这正是第~\ref{sec:pseudo} 节"LLM 含两个事件句 = 97.4\%"这一统计的直观来源。
规则法在该样本上给出一致答案，也是"精确率 0.985"的具体体现。

\subsection{实验命令与日志对照}

\begin{table}[H]
\centering
\caption{实验与其启动参数、日志文件}
\label{tab:runs}
\small
\begin{tabular}{lll}
\toprule
实验 & 关键参数 & 日志\\
\midrule
E1 完整方法 & \texttt{--cgl\_w 6.0} & \texttt{logs\_E1\_full.txt}\\
E0 去 CGL & \texttt{--cgl\_w 0} & \texttt{logs\_E0\_no\_cgl.txt}\\
E2 弱 CGL & \texttt{--cgl\_w 3.0} & \texttt{logs\_E2\_cgl3.txt}\\
快速检查 & \texttt{--k\_fold 2 --epoch 1} & \texttt{logout/seed\_209/}\\
\bottomrule
\end{tabular}
\end{table}

\subsection{本文对原始代码的全部改动}

为便于审阅，列出所有改动点。除标注"改动数值"者外，均不改变模型行为。

\begin{table}[H]
\centering
\caption{对官方仓库的改动清单}
\label{tab:changes}
\small
\begin{tabularx}{\textwidth}{@{}l l X@{}}
\toprule
文件 & 类型 & 内容\\
\midrule
\texttt{model/model.py} & 性能 & 5 处向量化改写（表~\ref{tab:patch}），已验证数值等价\\
\texttt{model/model.py} & 功能 & CGL 权重 \texttt{w=6} 改为读取 \texttt{args.cgl\_w}（默认 6.0）\\
\texttt{model/model.py} & 修复 & 删除死代码 \texttt{true\_sent}；修复 \texttt{sent\_ground} 重复追加\\
\texttt{model/\_orig\_model.py} & 新增 & 原版备份，供等价性验证\\
\texttt{parameter.py} & 功能 & 新增 \texttt{--grad\_ckpt}、\texttt{--cgl\_w}\\
\texttt{main.py} & 功能 & 每折按验证 $F_1$ 保存权重至 \texttt{checkpoints/<method>/}\\
\texttt{main.py} & 兼容 & 日志目录 \texttt{seed:\{s\}} $\to$ \texttt{seed\_\{s\}}（Windows 不允许冒号）\\
\texttt{utils.py} & 补齐 & 自 iLIF 取得\\
\texttt{tools/*} & 新增 & 9 个工具脚本（见第~\ref{sec:reproduce} 节目录树）\\
\bottomrule
\end{tabularx}
\end{table}
